\chapter[Linear and logistic models]{Linear and logistic models in high dimension, regularization}

The linear model is a fundamental and widely used approach in supervised
learning. It is based on the idea of exploiting a linear combination of explanatory
variables (\emph{features}) to generate predictions.

\section{Reminders concerning Linear Regression}

Linear regression assumes a linear relationship between the target variable
(\emph{label}) $Y$ and the explanatory variables $X$. Although this assumption is often
a simplified approximation, it proves extremely useful both conceptually and
practically.

\subsection{Formulation of the Linear Regression Model}

Consider a dataset following the hypothesis in Section~2.1.3 denoted
$\cD_N = \{(X_i, Y_i),\ i = 1, \ldots, N\}$. Here, $X$ is a variable taking values in $\R^p$,
representing the $p$ explanatory variables (\emph{features}), and $Y$ is a variable taking
values in $\R$, corresponding to the target variable (or dependent variable). For all
$i \in \{1, \ldots, N\}$, we have $X_i = (x_{i,1}, \ldots, x_{i,p}) \in \R^p$ and, to simplify
notations, we assume that $x_{i,1} = 1$\footnote{This assumption allows to treat a possible term that you may have encountered under the name of ``bias'' or ``intercept''; with this convention the term enters nicely into the vectorial format; see section~4.2 for an example of intercept (the `$b$' constant).}.

\begin{tcolorbox}[remarkbox]
\textbf{Assumption (H$_1$):} The relationship between $Y$ and $X$ is assumed to be
linear, i.e., there exists a vector $\beta = (\beta_1, \ldots, \beta_p)\T \in \R^p$ such that:
\begin{equation}\label{eq:ch4-linear-model}
Y = X\beta + \varepsilon.
\end{equation}
\end{tcolorbox}

\begin{remark}
In Equation~\eqref{eq:ch4-linear-model}, $\beta = (\beta_1, \ldots, \beta_p)\T \in \R^p$ represents the
unknown parameters of the model, which we aim to estimate (the optimization
will be detailed in Sections~\ref{sec:ch4-ols} and~\ref{sec:ch4-mle}). The term $\varepsilon$ represents the
noise, a random variable on which additional assumptions will be formulated
later, see Assumption (H$_2$).
\end{remark}

Under Assumption (H$_1$), the data in $\cD_N$ can be written as follows:
\begin{equation}\label{eq:ch4-data-scalar}
Y_i = \beta_1 x_{i,1} + \beta_2 x_{i,2} + \cdots + \beta_p x_{i,p} + \varepsilon_i, \qquad i = 1, \ldots, N,
\end{equation}
where $\varepsilon_i \sim$ i.i.d.~$\varepsilon$. This expression can be reformulated in a matrix form:
\begin{equation}\label{eq:ch4-data-matrix}
\Y = \X \beta + \varepsilon,
\end{equation}
with:
\[
\Y = \begin{pmatrix} Y_1 \\ \vdots \\ Y_N \end{pmatrix}_{N \times 1}, \quad
\X = \begin{pmatrix} x_{1,1} & \cdots & x_{1,p} \\ \vdots & \ddots & \vdots \\ x_{N,1} & \cdots & x_{N,p} \end{pmatrix}_{N \times p}, \quad
\varepsilon = \begin{pmatrix} \varepsilon_1 \\ \vdots \\ \varepsilon_N \end{pmatrix}_{N \times 1}.
\]

\begin{tcolorbox}[remarkbox]
\textbf{Assumption (H$_2$):} The following assumptions are made:
\begin{enumerate}
\item $\rank(\X) = p$,
\item $\varepsilon \sim \Normal(0_N, \sigma^2 \Id_N)$, which means $\varepsilon_1, \ldots, \varepsilon_N \iid \Normal(0, \sigma^2)$.
\end{enumerate}
\end{tcolorbox}

Here $0_N$ and $\Id_N$ are the null $N \times N$ matrix and the identity $N \times N$ matrix
respectively. The goal is to determine an optimal estimator\footnote{``Optimal'' can be defined according to various criteria, such as the \emph{least squares} method (Section~\ref{sec:ch4-ols}), the \emph{maximum likelihood} method (Section~\ref{sec:ch4-mle}), or methods like \emph{Ridge} and \emph{Lasso} (Section~\ref{sec:ch4-lasso}).}
$\beta^\star = (\beta_1^\star, \ldots, \beta_p^\star)\T \in \R^p$ based on the data $\cD_N$, and construct a predictive model (linear regression):
\begin{equation}\label{eq:ch4-predictor}
\widehat{f}^{\cD_N}(x) = x \beta^\star.
\end{equation}
For a new data point $x_0 \in \R^p$, the prediction of this model will be $y_0 = x_0 \beta^\star$
($x_0$ is considered as a row vector).

\subsection{Overview of the Ordinary Least Squares Estimator}\label{sec:ch4-ols}

\begin{definition}[Ordinary Least Squares Estimator]\label{def:ch4-ols}
The (ordinary) least squares estimator $\widehat{\beta}^{(OLS)}$ is defined as:
\begin{equation}\label{eq:ch4-ols-def}
\widehat{\beta}^{(OLS)} \in \argmin_{\alpha \in \R^p} \norm{\Y - \X \alpha}^2
= \argmin_{\alpha \in \R^p} \sum_{i=1}^N \bigl(Y_i - X_i \alpha\bigr)^2,
\end{equation}
where $\norm{\cdot}$ denotes the Euclidean norm in $\R^n$.
\end{definition}

\begin{remark}
By Definition~\ref{def:ch4-ols}, the model $\widehat{f}^{(OLS)}(x) := x \widehat{\beta}^{(OLS)}$ is the
\textbf{linear} model that minimizes the empirical risk $\cR^{\cD_N, l_{MSE}}(\widehat{f})$ on $\cD_N$ for the
loss function $l_{MSE}$ in~(2.5).
However, this approach does not automatically address the problem of over-fitting.
Consequently, the model $\widehat{f}^{(OLS)}(x) := x \widehat{\beta}^{(OLS)}$ may result in over-fitting.
\end{remark}

\begin{proposition}\label{prop:ch4-ols-formula}
Under Assumptions (H$_1$) and (H$_2$), the least squares estimator for the linear regression is:
\begin{equation}\label{eq:ch4-ols-formula}
\widehat{\beta}^{(OLS)} = (\X\T \X)^{-1} \X\T \Y.
\end{equation}
\end{proposition}

\begin{proof}
This result is admitted and can be found in many statistical textbooks. It can be
obtained from the convexity of the loss and the first order critical point equations.
\end{proof}

\begin{remark}
To obtain additional properties of $\widehat{\beta}^{(OLS)}$ such as unbiasedness,
some hypotheses on $\varepsilon$ are required, for instance $\E[\varepsilon \mid X] = 0$.
\end{remark}

Under Assumptions (H$_1$) and (H$_2$), the model $\widehat{f}^{(OLS)} = x \widehat{\beta}^{(OLS)}$ is called
\emph{optimal} in the least squares sense.

\subsection{Overview of the Maximum Likelihood Estimator (MLE)}\label{sec:ch4-mle}

We first define the Maximum Likelihood Estimator for a general parametric
procedure and then we apply to our linear model. Let $Z$ be a real-valued
random variable, and suppose that the distribution of $Z$, denoted $\mu_Z^\theta$, depends
on a parameter $\theta \in \Theta$, where $\Theta$ is the set of all possible parameters. For
example:
\begin{itemize}
  \item If $Z$ has a discrete distribution, we denote $p_\theta(y) = \PP(Z = z)$.
  \item If $Z$ has a continuous distribution (with respect to the Lebesgue measure),
  we denote $\mu_Z^\theta(dz) = f_\theta(z) \lambda(dz)$, where $f_\theta(z)$ is the density of $Z$ with respect
  to the Lebesgue measure $\lambda(dz)$.
\end{itemize}

Let $\mathbf{z} = (z_1, \ldots, z_n)$ be independent and identically distributed (i.i.d.) realizations
of $Z$. The likelihood function\footnote{The likelihood can be interpreted as ``what is the probability of observing this sample $(z_1, \ldots, z_n)$ given the parameter $\theta$?''.} for $(z_1, \ldots, z_n)$ is given by:
\begin{align}
&\text{in the discrete case } L(\mathbf{z}; \theta) = L(z_1, \ldots, z_n; \theta) := \prod_{i=1}^n p_\theta(z_i),\label{eq:ch4-lik-disc}\\
&\text{in the continuous case } L(\mathbf{z}; \theta) = L(z_1, \ldots, z_n; \theta) := \prod_{i=1}^n f_\theta(z_i).\label{eq:ch4-lik-cont}
\end{align}

This expression may be difficult to treat algebraically and for this reason
it is often replaced by the log-likelihood function, which simplifies the computation
by converting products into sums; the log-likelihood is defined as:
\begin{equation}\label{eq:ch4-loglik-def}
\text{Log likelihood: } \ell(\mathbf{z}; \theta) = \log L(\mathbf{z}; \theta).
\end{equation}

\begin{definition}[Maximum Likelihood Estimator (MLE)]
If there exists a $\widehat{\theta} \in \Theta$ such that
\begin{equation}\label{eq:ch4-mle-def}
L(\mathbf{z}; \widehat{\theta}) = \sup_{\theta \in \Theta} L(\mathbf{z}; \theta) \quad \text{or equivalently: } \ell(\mathbf{z}; \widehat{\theta}) = \sup_{\theta \in \Theta} \ell(\mathbf{z}; \theta),
\end{equation}
then $\widehat{\theta}$ is called a\footnote{Note that the MLE is not necessarily unique.} Maximum Likelihood Estimator (MLE) and is often denoted
by $\widehat{\theta}^{(ML)}$\footnote{The idea of the Maximum Likelihood Estimator is to find a parameter $\theta$ such that the likelihood of observing the realizations $(y_1, \ldots, y_n)$ is maximized.}.
\end{definition}

We now specialize to the Maximum Likelihood Estimator for a linear regression.
The hypotheses (H$_1$) and (H$_2$) imply that $Y \sim \Normal(X\beta, \sigma^2)$, which is a distribution
(conditional on $X$) that depends on the parameter $\beta$. Consequently,
the likelihood and log-likelihood functions for the data $\cD_N = \{(X_i, Y_i),\ 1 \le i \le N\}$
are\footnote{Cf. appendix A.4.1}:
\begin{align}
L_{\cD_N}(\beta) &= \frac{1}{(2\pi \sigma^2)^{N/2}} \exp\!\left[ -\frac{1}{2\sigma^2} \sum_{i=1}^N \bigl(Y_i - X_i \beta\bigr)^2 \right], \nonumber\\
\ell_{\cD_N}(\beta) &= -\frac{N}{2} \log(2\pi) - \frac{N}{2} \log \sigma^2 - \frac{1}{2\sigma^2} \sum_{i=1}^N \bigl(Y_i - X_i \beta\bigr)^2. \label{eq:ch4-loglik-linreg}
\end{align}

According to~\eqref{eq:ch4-loglik-linreg}, it is clear that the estimator $\widehat{\beta}$ that maximizes this log-likelihood
is the estimator that minimizes $\sum_{i=1}^N (Y_i - X_i \beta)^2$, which corresponds
to the least squares estimator.

\begin{proposition}[MLE for linear regression]
Under the hypotheses (H$_1$) and (H$_2$), the maximum likelihood estimator of the linear model~\eqref{eq:ch4-data-scalar} is
\begin{equation}\label{eq:ch4-mle-equals-ols}
\widehat{\beta}^{(MLE)} = (\X\T \X)^{-1} \X\T \Y = \widehat{\beta}^{(OLS)}.
\end{equation}
\end{proposition}

\section{Logistic Regression for Binary Classification}\label{sec:ch4-logistic}

Here $\cY = \{0, 1\}$. We still denote $\cD_N = \{(X_i, Y_i),\ 1 \le i \le N\}$ the training
data for the model. Let $x = (x^1, \ldots, x^p)\T \in \cX = \R^p$, where $\cX$ represents the
space of $p$ features\footnote{Note that now inputs are represented as a column vector.}. A linear combination of $(x^1, \ldots, x^p, 1)\T$ can be expressed as:
\begin{equation}\label{eq:ch4-linear-comb}
w_1 x^1 + \cdots + w_p x^p + b = w\T x + b,
\end{equation}
where $w = (w_1, \ldots, w_p)\T \in \R^p$ is the weight vector (slope vector), and $b \in \R$
is the intercept (sometimes called ``bias'').

\begin{tcolorbox}[remarkbox]
\textbf{Intuition 4.8.} The question is whether the linear combination of $w\T x + b$
can be used to solve a binary classification problem i.e., with $\cY = \{0, 1\}$. A naive answer is negative because for a binary
problem, $\cY = \{0, 1\}$, whereas $w\T x + b$ takes values in $\R$.
Nevertheless, the linear model can be useful in the context of binary
classification if we introduce an activation function $\boldsymbol{\sigma} : \R \to [0, 1]$; this
activation $\boldsymbol{\sigma}(w\T x + b)$ will be used to predict the label conditional to
$X = x$: if $\boldsymbol{\sigma}(w\T x + b) > 0.5$, predict label 1; otherwise, predict label 0.
\end{tcolorbox}

\subsection{Definition of the Logistic Regression}

We introduce a particular activation function that will prove very useful in the
following.

\begin{definition}[Sigmoid activation function]
The Sigmoid activation function is defined as:
\begin{equation}\label{eq:ch4-sigmoid-def}
\boldsymbol{\sigma} : x \in \R \mapsto \boldsymbol{\sigma}(x) = \frac{1}{1 + e^{-x}} \in\; ]0, 1[.
\end{equation}
\end{definition}

\begin{lemma}
The sigmoid function has the following properties:
\begin{align}
\forall x \in \R:\quad & 0 < \boldsymbol{\sigma}(x) < 1 \label{eq:ch4-sig-bounds}\\
& \boldsymbol{\sigma}(x)' = \boldsymbol{\sigma}(x)\bigl(1 - \boldsymbol{\sigma}(x)\bigr) \label{eq:ch4-sig-deriv}\\
& \boldsymbol{\sigma}(-x) + \boldsymbol{\sigma}(x) = 1. \label{eq:ch4-sig-sym}\\
& \boldsymbol{\sigma}^{-1}(\cdot) = \mathrm{logit}(\cdot),\ \text{where } \mathrm{logit}(p) = \log\!\left( \frac{p}{1-p} \right). \label{eq:ch4-sig-logit}
\end{align}
\end{lemma}
\begin{proof}
By direct computations.
\end{proof}

Various generalizations can be proposed, for instance $\boldsymbol{\sigma}_{K,a,r}(t) = \frac{K}{1 + a e^{-rt}}$
with $K, a, r \in \R_+$. However, here we will use the standard Sigmoid function,
i.e., $K = a = r = 1$. See Figure~\ref{fig:ch4-sigmoid} for an illustration.
We move now to the definition of the logistic regression. In a logistic
regression model, the sigmoid function $\boldsymbol{\sigma}$ is used as an activation function.

\begin{figure}[h]\centering
% FIGURE PLACEHOLDER (uncomment & replace with \includegraphics when image is available):
% \fbox{\parbox{0.7\linewidth}{\centering\textit{[Figure from PDF page 45: plot of the sigmoid function $\mathrm{sigmoid}(x) = 1/(1+e^{-x})$ on $[-5, 5]$, with horizontal asymptote at $1$.]}}}
\caption{The sigmoid activation function.}\label{fig:ch4-sigmoid}
\end{figure}

The binary classifier defined by the logistic regression is:
\begin{equation}\label{eq:ch4-classifier}
\mathrm{classifier}^{\cD_N}(x) = \begin{cases} 1, & \text{if } \boldsymbol{\sigma}(w\T x + b) = \dfrac{1}{1 + \exp\!\bigl(-(w\T x + b)\bigr)} \ge 0.5 \\ 0, & \text{otherwise} \end{cases},
\end{equation}
where:
\begin{itemize}
\item $x \in \cX = \R^p$ is the input point to predict;
\item the weight vector $w = (w_1, \ldots, w_p)\T \in \R^p$ and the intercept $b \in \R$ are model parameters to be optimized using the data $\cD_N$;
\item $0.5$ is a decision threshold. This threshold can be adjusted to make the predictor more or less strict.
\end{itemize}

\begin{tcolorbox}[remarkbox]
\textbf{Intuition 4.11.} The main idea of $\{0, 1\}$-classification is to predict the
conditional probability $\PP(Y = 1 \mid X = x)$. For instance, in a $k$-NN
model, $\PP(Y = 1 \mid X = x)$ is predicted by $\frac{1}{K} \sum_{j=1}^K Y_{i_j(x)}$, which is the
average label of the $K$ nearest neighbors of $x$. A logistic regression
model is a model parametrized by $w, b$ where we suppose that
\begin{align}
&\PP(Y = 1 \mid X = x) = f^{(w,b)}(x) \label{eq:ch4-cond-prob-1}\\
&\text{with } f^{(w,b)}(x) := \boldsymbol{\sigma}(w\T x + b) = \frac{1}{1 + \exp\!\bigl(-(w\T x + b)\bigr)}. \label{eq:ch4-fwb-def}
\end{align}
The parameters $(w, b)$ will be searched using the training data $\cD_N = \{(X_i, Y_i),\ 1 \le i \le N\}$. This model formulation can seem unusual at
first but is just a way to specify parameters on some conditional distribution, here Bernoulli; compare with the linear model where assumptions
concerned parameters of another conditional distribution, a Gaussian.
\end{tcolorbox}

Other interpretations of the logistic regression exist.
\begin{enumerate}
\item Logistic regression can be viewed as linear regression of the \emph{logit} term.
Recall that the \emph{logit} defined in~\eqref{eq:ch4-sig-logit} that is, the logarithm of the odds,
maps probabilities in $(0, 1)$ onto $\R$. In our case, if, according to the
model, $\PP(Y = 1 \mid X = x) = \frac{1}{1 + \exp(-(w\T x + b))}$, then: $\PP(Y = 0 \mid X = x) =
1 - \PP(Y = 1 \mid X = x) = \frac{\exp(-(w\T x + b))}{1 + \exp(-(w\T x + b))}$. This implies: $\exp(w\T x + b) = \frac{\PP(Y = 1 \mid X = x)}{\PP(Y = 0 \mid X = x)}$ thus
\begin{equation}\label{eq:ch4-logit-regression}
w\T x + b = \log\!\left( \frac{\PP(Y = 1 \mid X = x)}{1 - \PP(Y = 1 \mid X = x)} \right) = \mathrm{logit}\bigl(\PP(Y = 1 \mid X = x)\bigr).
\end{equation}
So we are performing a regression on $\mathrm{logit}(\PP(Y = 1 \mid X = x))$ viewed as
depending on the value $x$. But, in this interpretation, we do not dispose
of stable empirical estimations of the logit.

\item Logistic regression can also be viewed as an artificial neural network with
no hidden layers and a single output, see Figure~\ref{fig:ch4-single-neuron} for an illustration
and the Chapter~7 for details on neural networks.
\end{enumerate}

\begin{figure}[h]\centering
% FIGURE PLACEHOLDER (uncomment & replace with \includegraphics when image is available):
% \fbox{\parbox{0.7\linewidth}{\centering\textit{[Figure from PDF page 46: a neural-network diagram showing inputs $x_1, x_2, x_3, \ldots, x_p$ each connected with weights $w_1, \ldots, w_p$ to a summation node $\sum_{i=1}^p w_i x_i + b$ (with bias $b$), whose output is passed through $\mathrm{sigmoid}(\cdot)$ to produce $\widehat{y}$.]}}}
\caption{Logistic Regression as a Single Neuron Network.}\label{fig:ch4-single-neuron}
\end{figure}

So far, we have defined the logistic regression predictor. The following
questions arise:
\begin{itemize}
\item How to find the parameters $w$ and $b$ by optimization?
\item In what sense is the optimization considered optimal?
\item What loss function should be chosen?
\end{itemize}

\subsection{Generalized Loss Function for Logistic Regression}

In Section~2, we introduced an example of a loss function for classification, the
binary loss in equation~(2.7). However, this function lacks desirable properties
(continuity, differentiability, convexity, etc.), making it challenging to apply
optimization techniques such as gradient descent to find the optimal values of
$(w, b)$. In this section, we define the loss function differently: \emph{via likelihood} or
\emph{via cross-entropy}. We work under assumptions~\eqref{eq:ch4-cond-prob-1} and~\eqref{eq:ch4-fwb-def}.

\subsubsection*{By Likelihood}

Under the logistic regression model:
\begin{equation}\label{eq:ch4-bern-params}
\PP(Y = 1 \mid X = x) = f^{(w,b)}(x), \qquad \PP(Y = 0 \mid X = x) = 1 - f^{(w,b)}(x).
\end{equation}
i.e., the conditional distribution of the variable $Y$ given $X = x$ follows a
Bernoulli distribution with parameter $f^{(w,b)}(x)$. Consequently for a dataset
$\cD_N$ the likelihood function can be expressed as:
\begin{align}
L_{\cD_N}(w, b) &= \prod_{1 \le i \le N} \PP(Y = Y_i \mid X = X_i) \nonumber\\
&= \prod_{1 \le i \le N\,:\, Y_i = 1} \PP(Y = 1 \mid X = X_i) \cdot \prod_{1 \le i \le N\,:\, Y_i = 0} \PP(Y = 0 \mid X = X_i) \label{eq:ch4-lik-split}\\
&= \prod_{1 \le i \le N\,:\, Y_i = 1} f^{(w,b)}(X_i) \cdot \prod_{1 \le i \le N\,:\, Y_i = 0} \bigl(1 - f^{(w,b)}(X_i)\bigr) \nonumber\\
&= \prod_{1 \le i \le N} f^{(w,b)}(X_i)^{Y_i} \bigl(1 - f^{(w,b)}(X_i)\bigr)^{1 - Y_i}. \label{eq:ch4-lik-bern}
\end{align}

The log-likelihood $\ell_{\cD_N}(w, b) = \log\bigl(L_{\cD_N}(w, b)\bigr)$ is thus:
\begin{equation}\label{eq:ch4-loglik-logistic}
\ell_{\cD_N}(w, b) = \sum_{1 \le i \le N} \Bigl( Y_i \log\bigl( f^{(w,b)}(X_i) \bigr) + (1 - Y_i) \log\bigl( 1 - f^{(w,b)}(X_i) \bigr) \Bigr).
\end{equation}

We aim to find a maximum likelihood estimator
\[
(w^*, b^*) \in \argmax L_{\cD_N}(w, b).
\]

Note that finding $\argmax L_{\cD_N}(w, b)$ is equivalent to finding $\argmax \ell_{\cD_N}(w, b)$,
which is also equivalent to finding $\argmin -L_{\cD_N}(w, b)$, and finally to finding
$\argmin -\ell_{\cD_N}(w, b)$. This allows us to define the cost function (commonly referred
to as the \emph{loss function} or \emph{cost function})\footnote{The terms ``loss function'' and ``cost function'' are analogous. The term ``loss function'' typically refers to the error for a single training data point, whereas ``cost function'' often refers to the average (or sum) of the loss functions over the entire training dataset.} coherent with the empirical risk
$\cR^{\cD_N}(f)$ defined in~(2.9), as:
\[
\cR^{\cD_N}\bigl( f^{(w,b)} \bigr) := -\ell_{\cD_N}(w, b),
\]
and we seek to find a minimizer of the function $(w, b) \mapsto -\ell_{\cD_N}(w, b)$.

\begin{remark}
The function $(w, b) \mapsto -\ell_{\cD_N}(w, b)$ is convex, cf. Exercise~4.2.
\end{remark}

\subsubsection*{By Cross-Entropy}

\begin{tcolorbox}[remarkbox]
\textbf{Intuition 4.13.} We can consider that the logistic regression outputs
not a number but a distribution; when the activated output is `$a$' we can
interpret it as a Bernoulli distribution with parameter `$a$'. On the other
hand, the dataset $\cD_N$ is also a collection of conditional distributions,
some may be Dirac masses (a particular case of Bernoulli distribution)
but other can be non-trivial Bernoulli distributions; this latter case
appears when in $\cD_N$ we have some $X_i = X_j$ but $Y_i \ne Y_j$. Once this
is done we need to operate with some (pseudo-) metric defined on the
space of Bernoulli distributions.
\end{tcolorbox}

We recall first a general definition of the cross-entropy.

\begin{definition}[Cross-entropy for two discrete probability measures]
Let $p, q$ be two probability measures defined on the same finite set $\mathcal{Z} = \{z_1, \ldots, z_K\}$.
The cross-entropy of $q$ relative to $p$ on $\mathcal{Z}$ is defined as:
\begin{equation}\label{eq:ch4-ce-discrete}
H(p, q) := -\sum_{z \in \mathcal{Z}} p(z) \log\bigl( q(z) \bigr).
\end{equation}
Here, $p$ is referred to as the reference probability measure.
\end{definition}

\begin{tcolorbox}[advancedbox]
\textbf{Advanced topics 4.15.} For two probability measures $p, q$ on an arbitrary space the entropy $H(p)$ \& cross-entropy $H(p, q)$ are defined as:
\begin{align}
\text{entropy:} \quad & H(p) = -\E_p[\log p], \label{eq:ch4-entropy-def}\\
\text{cross-entropy:} \quad & H(p, q) = -\E_p[\log q]. \label{eq:ch4-ce-def}
\end{align}
Then cross-entropy can be seen as the linearization in $q$ of the entropy
because it is a first order linear expansion at $q$: $H(p, q) = H(q) + \grad_q H(q) (p - q)$.
When measure $p$ has continuous density, the entropy $H(p)$ is also called
`differential entropy'. This case is more delicate because we are taking
the expectation of logarithm of the density; but density is with respect
to the Lebesgue measure which means that we have fixed the measurement
units, the (cross) entropy will then change if we change the units, so it
is not an intrinsic characteristic of the distribution.
\end{tcolorbox}

Cross-entropy is a concept from information theory that quantifies ``how
uncertain events are''. It can be understood here as a (non-symmetric) measure
of the difference between two probability distributions. In general, $H(p, q) \ne H(q, p)$,
so it is essential to specify which measure is the reference, here $p$.
Let us explain how cross-entropy applied in evaluating classification models.
For a given data point $x_0 \in \cX = \R^p$, if its true label is $y_0 = 1$ (we suppose
it is unique), it can be considered that we acquired some information on the
\emph{true} conditional probability $\PP[Y = y \mid X = x_0]$ of $Y$ given $X = x_0$, namely this
conditional distribution seems to equal the distribution $(1, 0)$. As a probability measure $p^{(x_0, y_0)}$ (conditioned on $X = x_0$) we obtain $p^{(x_0, y_0)}(1) = 1 = y_0$
and $p^{(x_0, y_0)}(0) = 0 = 1 - y_0$. On the other hand, when $y_0 = 0$ it is normal
to associate the distribution with probabilities $p^{(x_0, y_0)}(1) = 0 = y_0$ and
$p^{(x_0, y_0)}(0) = 1 = 1 - y_0$. So in both cases the associated distribution is
$(y_0, 1 - y_0)$.

Now, when $x_0$ is passed through the logistic regression model, we obtain
predicted conditional probabilities, denoted $q^{(x_0, y_0)}$, as follows:
\begin{align}
q^{(x_0, y_0)}(1) &= \PP(Y = 1 \mid X = x_0) = f^{(w,b)}(x_0), \label{eq:ch4-q-one}\\
q^{(x_0, y_0)}(0) &= \PP(Y = 0 \mid X = x_0) = 1 - f^{(w,b)}(x_0). \label{eq:ch4-q-zero}
\end{align}

\begin{remark}
The sigmoid $\boldsymbol{\sigma}(\cdot)$ does not take values $0$ or $1$ so the predicted
probabilities are never exactly $0$ or $1$.
\end{remark}

Using $p^{(x_0, y_0)}$ as the reference measure in the cross-entropy definition, we
compute the cross-entropy of the predicted distribution relative to the conditional probability for this data point $(x_0, y_0)$ as:
\begin{align}
H\bigl(p^{(x_0, y_0)}, q^{(x_0, y_0)}\bigr) &= -\sum_{y \in \cY} p^{(x_0, y_0)}(y) \log\bigl( q^{(x_0, y_0)}(y) \bigr) \nonumber\\
&= -y_0 \cdot \log\bigl( q^{(x_0, y_0)}(y) \bigr) - (1 - y_0) \cdot \log\bigl( q^{(x_0, y_0)}(y) \bigr) \nonumber\\
&= -\Bigl[ y_0 \cdot \log\bigl( f^{(w,b)}(x_0) \bigr) + (1 - y_0) \cdot \log\bigl( 1 - f^{(w,b)}(x_0) \bigr) \Bigr]. \label{eq:ch4-ce-pointwise}
\end{align}

This leads to consider for the training data $\cD_N$, a cost function based on
cross-entropy defined as:
\begin{align}
\cR^{\cD_N}\bigl( f^{(w,b)} \bigr) &:= \frac{1}{N} \sum_{i=1}^N H\bigl( p^{(X_i, Y_i)}, q^{(X_i, Y_i)} \bigr) \nonumber\\
&= \frac{1}{N} \sum_{i=1}^N -\Bigl[ Y_i \cdot \log\bigl( f^{(w,b)}(X_i) \bigr) + (1 - Y_i) \cdot \log\bigl( 1 - f^{(w,b)}(X_i) \bigr) \Bigr] \nonumber\\
&= -\frac{1}{N} \ell_{\cD_N}(w, b), \label{eq:ch4-ce-cost}
\end{align}
where $\ell_{\cD_N}(w, b)$ is the log-likelihood from~\eqref{eq:ch4-loglik-logistic}. Note that strictly speaking
this is \textbf{not} a loss function in the sense of Definition~2.2 but it still has the same
intention. So we proved that
\begin{equation}\label{eq:ch4-mle-ce-equiv}
\begin{gathered}
\text{\textbf{maximizing the (log-)likelihood (averaged over the dataset)}} \\
\text{is equivalent to} \\
\text{\textbf{minimizing the cross-entropy (averaged over the dataset)}}.
\end{gathered}
\end{equation}

\begin{tcolorbox}[advancedbox]
\textbf{Advanced topics 4.17.} One may ask what happens when the empirical
distribution of $Y$ conditional to $X = x$ is not a Dirac mass but a general
Bernoulli distribution. One has to prove that the log-likelihood / entropy
is additive with respect to a mixture of distributions. This seems obvious
when the mixture concerns the $p$ distribution but not $q$, which is enough
for our purpose. See Exercise~4.3 for a more formal definition. This
may seem at first contradictory with the optimization objective: for a
data point $x_0, y_0$ our loss tends to orient $q^{x_0, \cdot}$ towards the Dirac mass in
$y_0$ while it also orients $q^{x_0, \cdot}$ towards $y_1$ should $(x_0, y_1)$ also appear in the
dataset. This is because $q^{\cdot, \cdot}$ truly depends only on the first argument.
So we never really expect the loss to be completely null unless there is
deterministic dependence.
\end{tcolorbox}

\subsection{Finding the solution}

Contrary to the situation of the linear regression, the minimization above does
not have a close formula for the optimal solution. Numerical algorithms are
then used to find it. It has been proved, cf.~\cite{turinici1} that for non-pathological data\footnote{Data is pathological if for instance all labels are equal or the input vectors are collinear.}
the properties of the optimal point depend on whether the dataset is separable
by a hyperplane or not; the surprising result is that the separable case is not
very convenient because in this case the coefficients of the regression go to
infinity, see Example~\ref{ex:ch4-separable} and Exercise~4.4.

\begin{example}\label{ex:ch4-separable}
Consider a simple example with the dataset $\cD = \{(x_1, y_1), (x_2, y_2)\}$
with all $x_i, y_i$ real numbers. The log-likelihood for the logistic regression with parameters $(w, b)$ is $\ell_{\cD}(w, b) = \sum_{i=1}^2 \bigl[ y_i \log \boldsymbol{\sigma}(x_i w + b) + (1 - y_i) \log\bigl( 1 - \boldsymbol{\sigma}(x_i w + b) \bigr) \bigr]$ where $\boldsymbol{\sigma}(\cdot)$ is the sigmoid activation defined in~\eqref{eq:ch4-sigmoid-def}.
Take the \textbf{separated} dataset with $x_1 = -1$, $y_1 = 0$ and $x_2 = 1$, $y_2 = 1$, the log-likelihood becomes: $\widetilde{\ell}_{\cD}(w, b) = \log\bigl( 1 - \boldsymbol{\sigma}(-w + b) \bigr) + \log \boldsymbol{\sigma}(w + b)$. Denoting $\alpha = e^w$, $\beta = e^b$ we obtain the function to maximize as $\log\!\left( \frac{\alpha^2 \beta}{(\alpha + \beta)(1 + \alpha \beta)} \right)$ which,
since $\alpha, \beta \ge 0$, is maximized for $\alpha = \infty$ and $\beta$ any finite value. But $\alpha = \infty$
corresponds to $w = \infty$!
\end{example}

Since analytic solutions are not available we will have to resort to numerical
procedures. This will be discussed in Sections~\ref{sec:ch4-ridge} and~\ref{sec:ch4-lasso} and Chapter~5.

\subsection{The Decision Boundary of the Logistic Regression}

The exact decision boundary of a $\{0, 1\}$-classification model is the set:
\begin{align}
\text{Exact Decision Boundary} &= \bigl\{ x \in \cX = \R^p \mid \text{prediction of } \PP(Y = 1 \mid X = x) \nonumber\\
&\qquad\qquad\qquad\qquad = \text{prediction of } \PP(Y = 0 \mid X = x) \bigr\} \nonumber\\
&= \bigl\{ x \in \cX \mid f^{(w,b)}(x) = 1 - f^{(w,b)}(x) \bigr\} = \bigl\{ x \in \cX \mid f^{(w,b)}(x) = 1/2 \bigr\}. \label{eq:ch4-decision-boundary}
\end{align}
where we used that in the logistic regression model, the prediction of $\PP(Y = 1 \mid X = x)$
is estimated by $f^{(w,b)}(x)$ and the prediction for $\PP(Y = 0 \mid X = x)$
is estimated by $1 - f^{(w,b)}(x)$. Since $f^{(w,b)}(x) = \boldsymbol{\sigma}(w\T x + b) = \frac{1}{1 + \exp(-(w\T x + b))}$
the condition $f^{(w,b)}(x) = 1/2$ becomes $\exp(-(w\T x + b)) = 1$ which happens
when $w\T x + b = 0$ consequently, the decision boundary of a logistic regression
model is the hyperplane:
\begin{equation}\label{eq:ch4-hyperplane}
\bigl\{ x \in \cX \mid w\T x + b = 0 \bigr\}.
\end{equation}

So the logistic regression will be successful when the data is \textbf{linearly separated}.
What if this is not the case? One may for instance change variables to reduce
to the previous linear separation case, see Figure~\ref{fig:ch4-decision-bound}.

\begin{remark}
In reality we should distinguish between the exact and numerical decision boundary. The latter is similar to the exact version but
uses $\widehat{w}, \widehat{b}$ the logistic regression solutions obtained by previous steps.
\end{remark}

\begin{figure}[h]\centering
% FIGURE PLACEHOLDER (uncomment & replace with \includegraphics when image is available):
% \fbox{\parbox{0.9\linewidth}{\centering\textit{[Figure from PDF page 53: two scatter plots. Left: points in the $(x_1, x_2)$-plane with a linear (dashed) decision boundary separating two coloured classes. Right: points arranged so that one class forms an inner disk surrounded by the other class; the decision boundary is a dashed circle, suggesting a switch to polar coordinates to make the data linearly separable.]}}}
\caption{The logistic regression method is better suited for linearly separable data (left). One could switch to polar coordinates on the data on the right to make them linearly separable.}\label{fig:ch4-decision-bound}
\end{figure}

\section{Why regularization?}

We will give several examples of problems encountered in regression fitting.
As we will see, all are somehow related and will lead us to envision a common
solution.

\subsection{Data dependence}

Consider a simple 2D linear regression problem with four data points as in
Figure~\ref{fig:ch4-data-dep}.

We fit a simple linear regression model using three points as training set
and one for testing. Consider two situations: fitting on first three points $(1, 2)$,
$(2, 3)$, $(3, 5)$ and fitting on Points $(2, 3)$, $(3, 5)$, $(4, 10)$. We see in Figure~\ref{fig:ch4-data-dep} that
the two regression give completely different results for the points $X = 1$ and
$X = 4$. Depending on which points we choose for training, the resulting model
can vary significantly. Both models fit their respective training data well but
perform poorly on the test points. Which one is the best? How can we not
let our procedure be affected too much by a small modification in the data?

\begin{figure}[h]\centering
% FIGURE PLACEHOLDER (uncomment & replace with \includegraphics when image is available):
% \fbox{\parbox{0.9\linewidth}{\centering\textit{[Figure from PDF page 54: on the left, a table of four data points $(X, Y) \in \{(1, 2), (2, 3), (3, 5), (4, 10)\}$; on the right, a scatter plot of these four points together with two linear fits --- Model 1 trained on the first three points and Model 2 trained on the last three points --- showing very different predictions at $X = 1$ and $X = 4$.]}}}
\caption{Data dependence of the linear regression prediction. \textbf{Left:} the dataset. \textbf{Right:} two trained models using different training sets (subset of the initial dataset).}\label{fig:ch4-data-dep}
\end{figure}

\subsection{Data rank problems}

Another situation arrives when the dataset is high-dimensional and the number
of samples is insufficient to describe it completely; this can be seen in the rank
of the dataset matrix, see Figure~\ref{fig:ch4-pca-mnist}. In this case hypothesis H$_2$ is not satisfied
any more and the (numerical) solution may not be unique. Each solution will
produce a different prediction and we will witness overfitting.

\begin{figure}[h]\centering
% FIGURE PLACEHOLDER (uncomment & replace with \includegraphics when image is available):
% \fbox{\parbox{0.9\linewidth}{\centering\textit{[Figure from PDF page 54: two plots, one of the cumulative explained variance as a function of the number of principal components (rising sharply and levelling off below 1.0), and one of the eigenvalues (log scale) as a function of the principal-component index, showing a steep drop around index 713.]}}}
\caption{PCA of the MNIST dataset. At some point (index 713) the eigenvalues fall below numerical round-off threshold of $10^{-16}$ which means that the data is rank deficient.}\label{fig:ch4-pca-mnist}
\end{figure}

\subsection{Noise fitting, overfitting}

Finally, consider a third example, polynomial overfitting. Overfitting occurs
when a statistical model captures not only the underlying pattern of the data
but also the noise. This leads to excellent performance on the training data
but poor generalization to unseen data. Polynomial regression is a common
example where overfitting can be observed as the degree of the polynomial
increases. Polynomial regression is a simple linear regression but the input
data $X_i$ is augmented with all powers $X_i^k$ for $k$ from $0$ to some maximum
polynomial degree.

To demonstrate overfitting, we generated a noisy sine wave and fitted polynomials of varying degrees (up to $12$). The data was split into training and
testing sets to evaluate performance.

Figure~\ref{fig:ch4-poly-overfit} illustrates the fitted polynomials alongside the noisy data. As
the polynomial degree increases, the model fits the training data more closely,
but the error on the test data also increases, highlighting overfitting.

\begin{figure}[h]\centering
% FIGURE PLACEHOLDER (uncomment & replace with \includegraphics when image is available):
% \fbox{\parbox{0.9\linewidth}{\centering\textit{[Figure from PDF page 55: three plots. Top: the noisy training data together with fitted polynomials of degrees 3 and 12 (degree 3: train error 0.02, test error 0.04; degree 12: train error 0.00, test error 61208.55); the degree-12 curve oscillates wildly. Bottom left: train and test errors (log scale) as a function of polynomial degree, with the test error exploding at high degree. Bottom right: $L^2$ norm of the coefficient vector as a function of degree, increasing sharply with degree.]}}}
\caption{Polynomial regression fits with varying degrees. Training and test errors are displayed for each fit.}\label{fig:ch4-poly-overfit}
\end{figure}

For lower-degree polynomials (e.g., degree 1), the model underfits the data,
failing to capture the underlying sine wave pattern. As the degree increases,
the model captures more nuances of the training data, including noise, which
results in overfitting. This is evident from the increasing test error for higher-degree polynomials.
A closer look at the regression parameters shows that the coefficients for
the higher degree polynomials are very large. Large values will tend to be more
unstable. This leads to the regularization idea presented next.

\section{Ridge (L2) Regularization}\label{sec:ch4-ridge}

\begin{figure}[h]\centering
% FIGURE PLACEHOLDER (uncomment & replace with \includegraphics when image is available):
% \fbox{\parbox{0.9\linewidth}{\centering\textit{[Figure from PDF page 56: in the $(\beta_1, \beta_2)$-plane, a blue dot marks the OLS solution, concentric red dotted ellipses represent the contours of the residual sum of squares $\norm{\Y - \X \beta}_2^2$ centred on OLS, a blue circle represents the ridge penalty region $\norm{\beta}_2 \le M_\lambda^R$, and a black dot marks the ridge solution where the smallest ellipse touches the circle.]}}}
\caption{Axes $\beta_1$ and $\beta_2$ represent the coefficients in a 2D parameter space. Blue Dot (OLS): The point $(1, 1)$ is the unregularized least squares solution. Red dotted Ellipses: These are the residuals level curves centered at the OLS solution, i.e. the level lines of the data term $\norm{\Y - \X \beta}_2^2$. Blue Circle: The constraint region from the ridge penalty $\norm{\beta}_2 \le M_\lambda^R$. Black Dot (Ridge Solution): The point where the smallest ellipse (i.e., smallest data error) touches the blue constraint circle, representing the ridge solution.}\label{fig:ch4-ridge-geometry}
\end{figure}

\subsection{Ridge Regression (L2-regularization)}

We will explain the ridge regression idea for the linear regression. The principle
is to find a solution to the regression problem while ensuring its stability by
penalizing the L2 norm of the regression coefficients.

The most important properties follow from the proposition below:

\begin{proposition}\label{prop:ch4-ridge}
Let $\lambda \ge 0$ and define the penalized cost $\cR^{OLS, R, \lambda}(\beta)$:
\begin{equation}\label{eq:ch4-ridge-cost}
\cR^{OLS, R, \lambda}(\beta) := \norm{\Y - \X \beta}_2^2 + \lambda \norm{\beta}_2^2
= \sum_{i=1}^N (Y_i - X_i \beta)^2 + \lambda \beta_i^2
\end{equation}
where $\norm{\cdot}_2$ is the standard Euclidian norm defined in~(A.3). Also introduce
the matrix $T_\lambda$\footnote{This notation is related to Tichonov regularization, of which ridge is a particular case.}:
\begin{equation}\label{eq:ch4-T-lambda}
T_\lambda := \X\T \X + \lambda \Id_p,
\end{equation}
Then for $\lambda > 0$:
\begin{enumerate}
\item the optimization problem:
\begin{equation}\label{eq:ch4-ridge-min}
\min_{\beta \in \R^p} \cR^{OLS, R, \lambda}(\beta),
\end{equation}
has an unique solution
\begin{equation}\label{eq:ch4-ridge-sol}
\widehat{\beta}^{(R)} = T_\lambda^{-1} \X\T \Y,
\end{equation}
where $\Id_p = \left[\begin{smallmatrix} 1 & 0 & \cdots & 0 \\ 0 & 1 & \cdots & 0 \\ \vdots & \vdots & \ddots & \vdots \\ 0 & 0 & \cdots & 1 \end{smallmatrix}\right] \in \R^{p \times p}$ is the $p \times p$ identity matrix.

\item Moreover $\widehat{\beta}^{(R)}$ is the unique solution to the constrained optimization
problem:
\begin{equation}\label{eq:ch4-ridge-constrained}
\min_{\beta \in \R^p} \norm{\Y - \X \beta}^2 \quad \text{subject to } \norm{\beta}_2 \le M_\lambda^R,
\end{equation}
where $M_\lambda^R = \norm{T_\lambda^{-1} \X\T \Y}_2$.
\end{enumerate}
\end{proposition}

\begin{proof}
The cost can be written in the form
\begin{equation}\label{eq:ch4-ridge-expanded}
\cR^{OLS, R, \lambda}(\beta) = \inner{T_\lambda \beta}{\beta} - 2 \inner{\X^T \Y}{\beta} + \norm{\Y}_2^2.
\end{equation}
When $\lambda > 0$, the matrix $T_\lambda$ is symmetric positive definite hence invertible.
We are thus minimizing a positive quadratic form $\cR^{OLS, R, \lambda}(\beta)$ which leads
to the existence and uniqueness of the minimizer.
As an alternative proof we can invoke Proposition~A.5 because the function
$\cR^{OLS, R, \lambda}(\beta)$:
\begin{itemize}
\item is obviously continuous
\item is coercive because of the term $\norm{\beta}_2^2$
\item is strictly convex being the sum of a strictly convex function $\norm{\beta}_2^2$ and a
convex function $\norm{\Y - \X \beta}^2$; the latter is convex because it is a composition
of a linear and a convex function, see Lemma~A.3.
\end{itemize}
The first order critical point equation (derivatives have to be null) are
$T_\lambda \widehat{\beta}^{(R)} = \X\T \Y$ and lead to the formula~\eqref{eq:ch4-ridge-sol}. Note that $\widehat{\beta}^{(R)}$ always
exists (even when $p > N$).
For the second part note that on one hand the problem~\eqref{eq:ch4-ridge-constrained} has obviously
a solution which is $\widehat{\beta}^{(R)}$. But this works also the other way around: any
other solution to~\eqref{eq:ch4-ridge-constrained} will be at least a good minimizer of $\cR^{OLS, R, \lambda}(\beta)$
as $\widehat{\beta}^{(R)}$ is, which, by the uniqueness of the minimizer of $\cR^{OLS, R, \lambda}(\beta)$ gives
the conclusion.
\end{proof}

\begin{tcolorbox}[advancedbox]
\textbf{Advanced topics 4.21.} The equivalence of the two minimization procedures proved in this proposition is also useful in another areas like
the mean-variance Markowitz portfolio optimization in finance, see~\cite{turinici4},
Chapters 3--5 and Appendix A].
\end{tcolorbox}

\begin{definition}[Ridge Estimator]
The ridge estimator is the estimator corresponding to the solution to the penalized cost minimization problem in Proposition~\ref{prop:ch4-ridge}, with $\lambda \ge 0$ being the hyper-parameter of the model\footnote{The parameter $\lambda$ can be set using \emph{cross-validation}, see Section~\ref{sec:ch4-cv}.}.
\end{definition}

Figure~\ref{fig:ch4-ridge-geometry} illustrates the constrained optimization problem~\eqref{eq:ch4-ridge-constrained} in dimension $p = 2$.

Note that when $\lambda = 0$, $\widehat{\beta}^{(R)}$ equals the Ordinary Least Squares (OLS)
estimator, provided $\X\T \X$ is invertible. If $\X\T \X$ is not invertible it can be
shown that a minimal norm solution is selected cf. Exercise~4.5. We prove
below that when $\lambda \to 0$ the optimum risk $\cR^{OLS, R, \lambda}$ will converge to the one
for the non-regularized problem.

\begin{lemma}
Denote $\cR^{OLS, R, \lambda, *} = \inf_\beta \cR^{OLS, R, \lambda}(\beta)$. Then:
\begin{equation}\label{eq:ch4-ridge-limit-zero}
\lim_{\lambda \to 0} \cR^{OLS, R, \lambda, *} = \min_\beta \norm{\Y - \X \beta}^2 = \cR^{OLS, R, 0, *}.
\end{equation}
\end{lemma}

\begin{proof}
Denote $\widehat{\beta}^{(R), \lambda}$ be the ridge estimator for $\lambda$. Then:
\begin{align}
\norm{\Y - \X \beta}^2 + \lambda \norm{\beta}^2 &= \cR^{OLS, R, \lambda}(\beta) \ge \cR^{OLS, R, \lambda, *} \nonumber\\
&= \norm{\Y - \X \widehat{\beta}^{(R), \lambda}}^2 + \lambda \norm{\widehat{\beta}^{(R), \lambda}}^2 \nonumber\\
&\ge \norm{\Y - \X \widehat{\beta}^{(R), \lambda}}^2 \ge \min_\beta \norm{\Y - \X \beta}^2 = \cR^{OLS, R, 0, *}. \label{eq:ch4-ridge-chain}
\end{align}
For $\lambda \to 0$ we obtain that any accumulation point of $\cR^{OLS, R, \lambda, *}$ is between
$\cR^{OLS, R, 0, *}$ and $\norm{\Y - \X \beta}^2$. Now minimize over $\beta$ to obtain the result.
\end{proof}

When $\lambda$ is large the answer is given by the next result.

\begin{lemma}
When $\lambda \to \infty$ the ridge estimator converges to the null vector.
\end{lemma}

\begin{proof}
Let $\widehat{\beta}^{(R), \lambda}$ be the ridge estimator for $\lambda$. By the minimization property
we can write
\begin{equation}\label{eq:ch4-ridge-bound-infty}
\lambda \norm{\widehat{\beta}^{(R), \lambda}}_2^2 \le \cR^{OLS, R, \lambda}(\widehat{\beta}^{(R), \lambda}) \le \cR^{OLS, R, \lambda}(0) = \norm{\Y}_2^2
\end{equation}
which implies $\norm{\widehat{\beta}^{(R), \lambda}}_2^2 \le \frac{\norm{\Y}_2^2}{\lambda}$ and we obtain the conclusion when $\lambda \to \infty$.
\end{proof}

\begin{remark}
It is possible to obtain an analytic formula for $\widehat{\beta}^{(R)}$, see Exercise~4.5.
\end{remark}

\begin{tcolorbox}[advancedbox]
\textbf{Advanced topics 4.26.} Some textbooks or software (e.g., the \texttt{sklearn} Python package) use another definition of the Ridge estimator as
the solution to the following penalized optimization problem:
\begin{equation}\label{eq:ch4-ridge-sklearn}
\min_{\beta \in \R^p} \norm{\Y - \X \beta}_2^2 + \lambda \sum_{j=1}^p \beta_j^2,
\end{equation}
where the difference is that intercept $\beta_0$ is not penalized. This comes
from the fact that in practice, before applying the ridge regression, data
is often standardized:
\begin{equation}\label{eq:ch4-standardize}
\begin{gathered}
X_{i,j} \to X_{i,j}^{(\mathrm{stand})} := \frac{X_{i,j} - \bar{X}_j}{\sigma_j}, \qquad Y_i \to Y_i^{(\mathrm{stand})} := \frac{Y_i - \bar{Y}}{\sigma_Y} \\
\text{with: } \bar{X}_j = \frac{\sum_{i=1}^N X_{i,j}}{N},\ \sigma_j = \left[ \frac{\sum_{i=1}^N (X_{i,j} - \bar{X}_j)^2}{N - 1} \right]^{1/2} \\
\bar{Y} = \frac{\sum_{i=1}^N Y_i}{N},\ \sigma_Y = \left[ \frac{\sum_{i=1}^N (Y_i - \bar{Y})^2}{N - 1} \right]^{1/2}
\end{gathered}
\end{equation}
In the \texttt{sklearn} package, the parameter \texttt{fit\_intercept = True} or \texttt{False}
can control which definition is considered.
\end{tcolorbox}

\section{Lasso (L1) regularization}\label{sec:ch4-lasso}

We will investigate now a second type of regularization, the Lasso (least absolute shrinkage and selection operator) regularization, also called $L^1$ regularization. It is introduced using an analogue of the Proposition~\ref{prop:ch4-ridge} that presents
its most important properties.

\begin{proposition}\label{prop:ch4-lasso}
Define the penalized cost:
\begin{equation}\label{eq:ch4-lasso-cost}
\cR^{OLS, L, \lambda}(\beta) := \norm{\Y - \X \beta}_2^2 + \lambda \norm{\beta}_1 = \sum_{i=1}^N (Y_i - X_i \beta)^2 + \lambda \abs{\beta_i}
\end{equation}
where $\norm{\cdot}_1$ is the L1 norm defined in~(A.2). Then
\begin{enumerate}
\item the optimization problem:
\begin{equation}\label{eq:ch4-lasso-min}
\min_{\beta \in \R^p} \cR^{OLS, L, \lambda}(\beta),
\end{equation}
has at least a solution $\widehat{\beta}^{(L)}$.

\item for any two solutions $\widehat{\beta}^{(L), 1}$ and $\widehat{\beta}^{(L), 2}$:
\begin{equation}\label{eq:ch4-lasso-same-pred}
\X \widehat{\beta}^{(L), 1} = \X \widehat{\beta}^{(L), 2} \quad \text{and} \quad \norm{\widehat{\beta}^{(L), 1}}_1 = \norm{\widehat{\beta}^{(L), 2}}_1.
\end{equation}

\item We will denote $M_\lambda^L := \norm{\widehat{\beta}^{(L)}}_1$ where $\widehat{\beta}^{(L)}$ is a solution. The problem~\eqref{eq:ch4-lasso-min} is equivalent to the constrained optimization problem:
\begin{equation}\label{eq:ch4-lasso-constrained}
\min_{\beta \in \R^p} \norm{\Y - \X \beta}^2 \quad \text{subject to } \norm{\beta}_1 \le M_\lambda^L.
\end{equation}

\item If $\rank(\X) = p$ then the solution is unique.
\end{enumerate}
\end{proposition}

\begin{proof}
We invoke Proposition~A.5 because the function $\cR^{OLS, L, \lambda}(\beta)$ is obviously continuous and coercive because of the term $\norm{\beta}_1$. Note that, as in
Proposition~\ref{prop:ch4-ridge} the function $\cR^{OLS, L, \lambda}(\beta)$ can be proven to be convex but,
unlike before, it is not necessarily strictly convex because here the norm
appears to power 1 not 2!
To prove~\eqref{eq:ch4-lasso-same-pred} note first that no risk can be below the risk of a minimizer
so the risk of all minimizers is the same. Also note that the function
\begin{equation}\label{eq:ch4-lasso-interp}
\xi \in [0, 1] \mapsto \norm{\Y - \X \bigl( (1 - \xi) \widehat{\beta}^{(L), 1} + \xi \widehat{\beta}^{(L), 2} \bigr)}^2
\end{equation}
is convex; if the function is strictly convex:
\begin{align}
\forall \xi \in\; ]0, 1[:\ \cR^{OLS, L, \lambda}\bigl( (1 - \xi) \widehat{\beta}^{(L), 1} + \xi \widehat{\beta}^{(L), 2} \bigr) &< (1 - \xi) \cR^{OLS, L, \lambda}(\widehat{\beta}^{(L), 1}) + \xi \cR^{OLS, L, \lambda}(\widehat{\beta}^{(L), 2}) \nonumber\\
&= \cR^{OLS, L, \lambda}(\widehat{\beta}^{(L), 1}) = \cR^{OLS, L, \lambda}(\widehat{\beta}^{(L), 2}), \label{eq:ch4-lasso-strict-conv}
\end{align}
which is impossible. But the function in~\eqref{eq:ch4-lasso-interp} is strictly convex as soon
as the coefficient $\norm{\X (\widehat{\beta}^{(L), 1} - \widehat{\beta}^{(L), 2})}_2^2$ of the quadratic term $\xi^2$ is non-null.
So we must have $\norm{\X (\widehat{\beta}^{(L), 1} - \widehat{\beta}^{(L), 2})}_2^2 = 0$ which gives the conclusion for
this part.
The proof of~\eqref{eq:ch4-lasso-constrained} is similar to the analogue statement in Proposition~\ref{prop:ch4-ridge}.
For the final conclusion note that when the rank is full we are back in the
situation of strict convexity so the uniqueness follows from Proposition~A.5.
\end{proof}

\begin{tcolorbox}[advancedbox]
\textbf{Advanced topics 4.28.} Using standard convexity arguments we can
even prove that either the solution is unique or there is an uncountable
set of solutions.
\end{tcolorbox}

\begin{definition}[Lasso Estimator]
The Lasso estimator is any solution of
the penalized cost minimization problem in Proposition~\ref{prop:ch4-lasso}, with $\lambda \ge 0$ being
the hyper-parameter of the model.
\end{definition}

Note that by Proposition~\ref{prop:ch4-lasso} point 2 all estimators share the same predictor so the LASSO predictor is well defined.

\begin{remark}
The Lasso regularization has a different homogeneity than
the ridge regularization, that is, the penalization term is of second order in
the ridge regularization and of first order in the lasso procedure. This will
have an impact on the behavior of the solution.
\end{remark}

Figure~\ref{fig:ch4-lasso-geometry} illustrates the constrained optimization problem~\eqref{eq:ch4-lasso-constrained} in dimension $p = 2$.

\begin{figure}[h]\centering
% FIGURE PLACEHOLDER (uncomment & replace with \includegraphics when image is available):
% \fbox{\parbox{0.9\linewidth}{\centering\textit{[Figure from PDF page 62: in the $(\beta_1, \beta_2)$-plane, a blue dot marks the OLS solution, concentric red dotted ellipses represent the residual-sum-of-squares contours, a rotated blue square (diamond) represents the Lasso penalty region $\norm{\beta}_1 \le M_\lambda^L$, and a black dot marks the Lasso solution at a corner of the diamond where the smallest ellipse touches the diamond.]}}}
\caption{Analog to Figure~\ref{fig:ch4-ridge-geometry} for Lasso regression.}\label{fig:ch4-lasso-geometry}
\end{figure}

What about extreme cases $\lambda = 0$ and $\lambda \to \infty$? Similar to the ridge
regression, when $\lambda = 0$ the Lasso estimator will equal the Ordinary Least
Squares (OLS) estimator, provided $\X\T \X$ is invertible; when $\lambda$ is large it will
converge to zero, as shown in the following lemma:

\begin{lemma}
When $\lambda \to \infty$ the Lasso estimators converge to the null vector.
\end{lemma}

\begin{proof}
Let $\widehat{\beta}^{(L), \lambda}$ be a Lasso estimator for $\lambda$. As in the proof of the
Lemma~4.24 we obtain by minimization:
\begin{equation}\label{eq:ch4-lasso-bound-infty}
\lambda \norm{\widehat{\beta}^{(L), \lambda}}_1 \le \cR^{OLS, L, \lambda}(\widehat{\beta}^{(L), \lambda}) \le \cR^{OLS, L, \lambda}(0) = \norm{\Y}_2^2
\end{equation}
which implies $\norm{\widehat{\beta}^{(L), \lambda}}_1 \le \frac{\norm{\Y}_2^2}{\lambda}$ and gives the conclusion when $\lambda \to \infty$.
\end{proof}

\begin{remark}
The Lasso regularization is mainly known for its preference
to \textbf{set to zero} some of the $\beta$ components. This allow to only keep relevant
explanatory variables and eliminate the others. Intuitively this behavior is
visible in the form of the level lines of the penalty, see Figure~\ref{fig:ch4-lasso-geometry}: while
for ridge regression level lines are curves, here level lines are unions of
linear objects and optimum is more likely to appear at the corners. See
Exercise~4.6 for technical details.
\end{remark}

\begin{tcolorbox}[advancedbox]
\textbf{Advanced topics 4.33.} Other regularization can be imagined combining ridge and lasso estimators, for instance Elasticnet will use a cost
term $\cR^{OLS, E, \lambda_1, \lambda_2}(\beta) := \norm{\Y - \X \beta}_2^2 + \lambda_1 \norm{\beta}_1 + \lambda_2 \norm{\beta}_2^2$.
\end{tcolorbox}

\begin{tcolorbox}[advancedbox]
\textbf{Advanced topics 4.34.} Regularization extends to any other procedures
that involve minimization, for instance to logistic regression.
\end{tcolorbox}

\section{Cross validation}\label{sec:ch4-cv}

Cross-validation is a robust statistical method used to estimate the performance of predictive models. Among various techniques, K-fold cross-validation
stands out due to its balance between bias and variance in model assessment.
We provide a formal description of the K-fold cross-validation procedure for a
general risk function and a dataset $\cD_N$.

Let $\cD_N = \{(X_i, Y_i)\}_{i=1}^N$ denote a dataset satisfying the hypothesis~(2.2)
where $X_i \in \cX$ represents the features and $Y_i \in \cY$ the corresponding response.
We consider a model $\widehat{f}(\cdot; \lambda)$ involving some hyper-parameter $\lambda \in \Lambda$ that aims
to approximate the true underlying relationship between $X$ and $Y$. The performance of the model is quantified by a risk function $\cR$ built from a loss
function $l$ as in~(2.9). The question is how to find the best hyper-parameter
of the model.

K-fold cross-validation involves partitioning the dataset $\cD_N$ into $K$ approximately equal-sized, mutually exclusive subsets (folds) $\cD_1, \cD_2, \ldots, \cD_K$.
The procedure is as follows:

\begin{enumerate}
\item \textbf{Partitioning the Data:} Divide $\cD_N$ into $K$ disjoint subsets $\cD_1, \cD_2, \ldots, \cD_K$ such that: $\bigcup_{k=1}^K \cD_k = \cD_N$, $\cD_k \cap \cD_l = \emptyset$ for $k \ne l$.

\item \textbf{Model Training and Validation:} For each fold $k \in \{1, 2, \ldots, K\}$:
\begin{enumerate}
\item Define the training dataset $\cD_{-k} = \cD_N \setminus \cD_k$.
\item Train the model on $\cD_{-k}$ to obtain $\widehat{f}^{(k)}(\cdot; \lambda)$;
\item Evaluate the model on the validation dataset $\cD_k$ using the risk
function $\cR$.
\end{enumerate}
\end{enumerate}

\begin{figure}[h]\centering
% FIGURE PLACEHOLDER (uncomment & replace with \includegraphics when image is available):
% \fbox{\parbox{0.9\linewidth}{\centering\textit{[Figure from PDF page 64: a grid illustrating 5-fold cross-validation: five rows labelled Fold 1, \ldots, Fold 5, each with five cells $\cD_1, \ldots, \cD_5$ marked either ``Val'' (validation) or ``Train'' (training); in each row a different single column is the validation set.]}}}
\caption{Illustration of K-Fold Cross-Validation with $K = 5$. Each fold is used once as a validation set while the remaining folds form the training set.}\label{fig:ch4-kfold}
\end{figure}

\begin{enumerate}
\setcounter{enumi}{2}
\item \textbf{Compute the Cross-Validation Estimate:} The K-fold cross-validation
estimate of the risk is given by $\widehat{\cR}_{\mathrm{CV}}(\lambda) = \frac{1}{K} \sum_{k=1}^K \cR^{\cD_k, l}\bigl( \widehat{f}^{(k)}(\cdot; \lambda) \bigr)$.

\item \textbf{Model Selection and Assessment:} The model with the lowest cross-validation risk estimate is selected, i.e., $\lambda_{\mathrm{opt}} \in \argmin \widehat{\cR}_{\mathrm{CV}}(\lambda)$.
\end{enumerate}

K-fold cross-validation strikes a balance between bias and variance:
\begin{itemize}
\item \textbf{Bias:} As $K$ increases, the training dataset in each fold becomes larger,
leading to a lower bias in the risk estimate.
\item \textbf{Variance:} A larger $K$ results in higher correlation between training sets
across folds, potentially increasing the variance of the risk estimate.
\end{itemize}

Common choices of $K$ include $K = 5$ or $K = 10$, which have been empirically shown to provide a good trade-off in many practical applications.

\section{Further techniques: Group Lasso}

Group Lasso is a regularization technique that extends the lasso to situations
where the coefficients are naturally organized into groups and one wants
to enforce that the whole group is either included in or excluded from the
model. Instead of applying the L1 penalty to individual coefficients (as in
standard lasso), group lasso applies the L2 penalty to each group of coefficients, but enforces the \textbf{L1 penalty at the group level}. This means that
either all coefficients in a group are included or none of them are included in
the final model.

\begin{sloppypar}
Consider a collection of coefficients $\beta_1, \beta_2, \ldots, \beta_p$ organized into $G$ groups
$\cG_1, \ldots, \cG_G$, $\bigcup_{g=1}^G \cG_g = \{1, \ldots, p\}$. The group lasso objective function can be
written as:
\end{sloppypar}
\begin{equation}\label{eq:ch4-group-lasso}
\text{Objective} = \text{Loss Function} + \lambda \sum_{g=1}^G \norm{\beta^g}_2,
\end{equation}
where:
\begin{equation}\label{eq:ch4-group-lasso-norm}
\begin{cases}
\norm{\beta^g}_2 = \sqrt{\sum_{i \in \cG_g} \beta_i^2}\ \text{is the L2 norm of the coefficients in group } \cG_g \\
\lambda > 0\ \text{is the regularization parameter}.
\end{cases}
\end{equation}

So, while \textbf{Lasso} applies the L1 penalty to individual coefficients, shrinking
some of them to zero, the \textbf{Group Lasso} applies the L2 penalty to the coefficients
within each group but applies the \textbf{L1 penalty} to the norms of these groups.
This means that the \textbf{entire group of coefficients} is either selected (if the
group norm is not shrunk to zero) or excluded (if the group norm is shrunk to
zero).

\begin{example}
Consider a regression model with several groups of features:
features $\beta_1, \beta_2$ related to the first category, features $\beta_3, \beta_4$ related to the second
category and feature $\beta_5$ related to the third category. With group lasso, the
model might choose to \textbf{keep both} $\beta_1$ and $\beta_2$ together (or discard them entirely)
but will not allow a situation where only one of the $\beta_1$ or $\beta_2$ is included while
the other is set to zero. The same applies to the other groups: either a group
is fully included, or it is excluded entirely.
\end{example}

This procedure can be applied to \textbf{categorical variables encoded as
multiple binary features}; in this case it is possible to select or exclude the
whole category at once. Also for \textbf{multivariate data} where multiple correlated features come from the same source (e.g., multiple sensors or repeated
\ldots)
